% 第2章 LaTeX 章节片段：文献综述。
% 引用条目位于 thesis_references.bib。

\chapter{人体上肢运动测量技术综述}
\label{chap:literature_review}

\section{研究范围与测量对象}
\label{sec:ch2_scope}

人体上肢运动测量的目标，是利用外部或可穿戴传感器，把连续人体动作转换为可计算的运动学
描述。根据应用目的，测量结果可以表示为人体关键点坐标、肢段位姿、肢段方向、关节角或手部
任务变量。设传感器原始观测为 $\boldsymbol y$，肢段状态为 $\boldsymbol x_s$，人体上肢广义
坐标为 $\boldsymbol q_h$，则一般测量过程可写为
\begin{equation}
  \boldsymbol y
  \xrightarrow{\mathcal P}
  \hat{\boldsymbol x}_s
  \xrightarrow{\mathcal K}
  \hat{\boldsymbol q}_h,
  \label{eq:ch2_measurement_chain}
\end{equation}
其中 $\mathcal P$ 表示从图像、惯性信号或电磁场观测恢复传感器/肢段状态，$\mathcal K$ 表示
依据人体运动学模型计算关节变量。两步之间并不总是能够严格分离：视觉方法可能直接输出关键点，
可穿戴方法则通常先输出传感器位姿，再通过安装关系恢复肢段方向。因而，传感器本身的测量精度
不能直接等同于人体关节角精度。

本文关注实时人机动作映射所需的左上肢测量。所需状态包括大臂方向、前臂方向、肩部三个旋转、
肘部屈伸以及前臂旋前/旋后。评价一种测量技术时，除静态精度外，还应考虑动态响应、轴向旋转
可观测性、遮挡或场畸变、佩戴复杂度、重新安装重复性、工作空间以及在线计算开销。本章据此
综述标记式光学动捕、无标记视觉、惯性/磁惯性测量、电磁跟踪和多模态参考方案。安装标定算法、
机器人动作重定向与安全执行、实验评价方法分别放在第4、5、6章结合本文方法进一步讨论。

\section{人体上肢运动测量方法}
\label{sec:ch2_human_motion_capture}

\subsection{上肢运动的表示层级}
\label{subsec:ch2_motion_representation}

上肢测量可以分为四个相互关联的表示层级。第一层是肩、肘、腕等关键点的二维或三维位置；
第二层是大臂、前臂等刚体肢段的位置和姿态；第三层是依据关节模型分解得到的肩、肘和腕关节角；
第四层是手部轨迹、抓取方向等任务变量。关键点位置便于视觉检测，但仅由肩、肘、腕三个点不能
唯一确定大臂和前臂绕自身纵轴的旋转。肢段六自由度位姿包含更多信息，却必须先定义传感器坐标
系与人体肢段坐标系之间的关系。关节角便于与机器人关节对应，但会受到旋转顺序、零位和人体
模型简化的影响。因此，比较不同系统时必须先确认其输出处于同一表示层级。

人体肩关节具有较大的运动范围，肩胛带运动还会造成关节中心相对躯干移动；肘部主要表现为屈伸，
前臂则可以绕纵轴旋前/旋后。若把上肢简化为由理想球铰肩、单轴肘和单轴前臂组成的串联链，模型
虽然便于实时计算，却不能完整描述肩胛骨、锁骨和腕部屈伸。本文所需的五自由度描述因此是一种
面向机器人控制的低维运动表示，而不是完整临床解剖模型。

\subsection{基于标记的光学运动捕捉}
\label{subsec:ch2_marker_mocap}

多相机光学运动捕捉通过三角测量恢复反光标记的三维坐标，并由同一肢段上的标记簇估计刚体
位姿。Cappozzo 等提出了骨骼嵌入解剖坐标系以及由体表标志点确定肢段位置和姿态的规范化
描述，为关节运动学的跨实验比较奠定了基础~\cite{cappozzo1995position}。该类系统具有较高的
空间分辨率和成熟的同步能力，因而常被用作可穿戴传感器的实验室参考。

但是，光学系统直接观测的是皮肤表面的标记，而不是骨骼。皮肤和肌肉相对骨骼的位移会产生
软组织伪影（soft tissue artefact, STA）；其大小随安装位置、动作幅度和受力状态变化，难以
通过一次静态偏置完全消除~\cite{leardini2005soft}。此外，多相机系统依赖视线，需要专用空间、
高质量标定和较复杂的数据处理。当手臂贴近躯干、旋转前臂或执行跨平面动作时，标记遮挡仍是
上肢测量的重要问题。

\subsection{基于普通相机的无标记视觉估计}
\label{subsec:ch2_markerless_vision}

无标记人体姿态估计从 RGB 或 RGB-D 图像中检测人体关键点。以 OpenPose 为代表的方法使用
部位亲和场同时估计关键点和肢段连接关系，实现了实时二维多人姿态估计
~\cite{cao2017realtime}。普通相机成本低、部署方便，也能提供直观的视频证据，适合大范围动作
和任务级演示记录。

其局限在于，单目图像存在深度歧义，二维关键点误差会在三维重建和关节角计算中放大；多相机
方法虽然改善深度估计，却引入内参、外参、同步和遮挡问题。学习模型的结果还受到训练数据、
衣着、光照和视角分布影响。对于本文关注的前臂旋前/旋后，仅依赖肘、腕二维关键点通常不能
稳定恢复肢段轴向旋转。因此，视觉更适合用作独立评价或辅助观测，而不是本研究实时控制链的
唯一输入。

\subsection{惯性与磁惯性可穿戴测量}
\label{subsec:ch2_imu}

惯性测量单元利用陀螺仪和加速度计估计姿态，磁惯性测量单元进一步使用磁力计提供航向参考。
Seel 等利用运动学约束识别关节轴和位置，说明在不预先精确对齐传感器外壳的情况下也可以
估计关节角~\cite{seel2014imu}。此类设备便携、工作空间大，且不依赖相机视线，已广泛用于
人体运动分析。

惯性方法仍需处理陀螺积分漂移、动态加速度和磁场扰动。更关键的是，传感器输出坐标系不等于
解剖肢段坐标系。Picerno 等通过解剖标志点建立上肢传感器到肢段的标定关系，并以光学系统
验证肩、肘运动学~\cite{picerno2019upper}；Zabat 等则表明多姿态标定可以减小安装偏差和软组织
运动造成的角度误差~\cite{zabat2019multiple}。近期研究进一步指出，当被测系统和参考系统采用
不一致的传感器到肢段标定定义时，即使原始姿态测量相近，所得关节角也可能出现明显系统偏置
~\cite{adjel2025influence}。这些结论对电磁可穿戴测量同样适用：传感器姿态准确并不自动等价于
解剖关节角准确。

\subsection{电磁位姿跟踪}
\label{subsec:ch2_em_tracking}

电磁跟踪器通过测量发射线圈与接收线圈之间的耦合，直接输出接收器相对发射器的三维位置和
三维姿态。与光学方法相比，它不要求传感器处于相机视野内，能够在人体自身遮挡时连续工作；
与纯惯性方法相比，它不依赖长时间积分得到位置。上述特性使其适合近距离、局部人体运动和
手术导航等缺少可靠视线的场景~\cite{maclachlan2025calibration}。

电磁系统的误差具有明显的空间依赖性。Milne 等发现，发射器与接收器距离和附近金属会影响
位置及方向精度~\cite{milne1996accuracy}。Frantz 等提出了区分准确度与精密度的标准化测试
协议~\cite{frantz2003accuracy}；Hummel 等采用该协议评价六自由度系统，进一步展示了工作空间
位置和金属材料对结果的影响~\cite{hummel2006evaluation}。MacLachlan 等从磁场模型标定和
表征角度说明，设备标定需要覆盖工作空间内的多个已知位姿，并且位置与方向误差存在耦合
~\cite{maclachlan2025calibration}。

因此，制造商给出的静态精度不能直接代表佩戴后的关节角精度。实际系统还会叠加接收器外壳与
肢段不平行、绑带滑动、软组织运动、无线数据丢失和两接收器不同步等误差。本文将发射器固定
在躯干上，以躯干为共同参考；两个接收器分别测量大臂和前臂。该布置减少了对外部世界坐标的
依赖，但仍必须在每次重新安装后估计传感器到肢段的旋转偏置和几何参数。

\subsection{多模态测量与独立视觉参考}
\label{subsec:ch2_multimodal_reference}

不同测量原理具有互补误差。视觉系统容易受到遮挡和深度歧义影响，但不会因电磁场畸变产生同类
误差；电磁系统不依赖视线，却对发射器距离和附近导体敏感。因此，将一种技术用于在线测量、
另一种物理原理独立的技术用于离线参考，比用同一传感链内部的拟合残差评价自身更有说服力。
多模态系统也可以进行在线融合，但这要求可靠的空间外参、时间同步和各模态不确定度模型。

低成本视觉参考通常包括相机标定和刚性视觉标记两部分。Zhang 的平面标定方法可由多个不同
姿态的平面图案估计相机内参和镜头畸变~\cite{zhang2000flexible}；AprilTag 和 ArUco 则通过
具有唯一编码的方形标记提供稳定角点和标记身份~\cite{olson2011apriltag,garrido2014aruco}。
两台固定相机共同观测安装于大臂和前臂的刚性标记时，可在统一视觉坐标系中重建肢段姿态，并与
电磁输出比较。该方案仍会受到打印尺寸、标记固定刚度、运动模糊和双相机同步误差影响，因此
应称为独立视觉参考，而不应未经不确定度分析直接称为绝对真值。

\subsection{测量方案比较与本文选择}
\label{subsec:ch2_sensing_comparison}

表~\ref{tab:ch2_sensing_comparison} 对与本研究相关的四类方案进行定性比较。不存在对所有指标
都占优的单一技术，选择取决于控制目标、工作空间和允许的实验复杂度。

\begin{table}[htbp]
  \centering
  \caption{人体上肢运动测量方案比较}
  \label{tab:ch2_sensing_comparison}
  \begin{tabularx}{\linewidth}{@{}lXXX@{}}
    \toprule
    方法 & 主要优点 & 主要限制 & 在本文中的角色 \\
    \midrule
    标记式光学动捕 & 多点三维测量成熟，便于统一实验室坐标 & 设备和场地成本高，存在遮挡与软组织伪影 & 文献中的常用参考方法 \\
    普通相机视觉 & 成本低，视频证据直观，非接触 & 深度歧义、遮挡、模型偏差，轴向旋转观测弱 & 双手机独立评价 \\
    IMU/MIMU & 便携、工作空间大、不依赖视线 & 漂移、动态加速度或磁干扰，需安装标定 & 技术路线对照 \\
    电磁跟踪 & 直接输出六自由度位姿，不受光学遮挡 & 工作范围有限，受金属和安装变化影响 & 实时控制主传感器 \\
    \bottomrule
  \end{tabularx}
\end{table}

\section{上肢测量中的共性问题}
\label{sec:ch2_common_challenges}

\subsection{空间可观测性与动作相关退化}
\label{subsec:ch2_observability}

测量精度并非在所有动作中恒定。单目视觉的深度方向、标记被遮挡的姿态、磁场工作空间边缘，
以及肢段方向接近共线的直臂姿态，都可能降低某些自由度的可观测性。尤其是前臂轴向旋转：仅靠
肩、肘、腕关键点不能恢复该旋转；即使传感器直接输出姿态，当肘部完全伸直时，也难以仅由两条
肢段方向区分肘角变化和轴向旋转。因此，实验动作必须覆盖多个平面和肘角，不能只用单一静态
姿势评价整个工作空间。

\subsection{坐标系、安装关系与重新佩戴}
\label{subsec:ch2_mounting_issue}

对于可穿戴传感器，传感器坐标系通常不与人体肢段解剖坐标系平行。设传感器坐标系为 $S$，
肢段坐标系为 $B$，则跟踪器报告的 ${}^{T}\boldsymbol R_S$ 必须结合安装变换
${}^{S}\boldsymbol R_B$ 才能得到肢段姿态：
\begin{equation}
  {}^{T}\boldsymbol R_B={}^{T}\boldsymbol R_S{}^{S}\boldsymbol R_B.
  \label{eq:ch2_sensor_segment_relation}
\end{equation}
重新佩戴会改变 ${}^{S}\boldsymbol R_B$，绑带滑动和软组织运动还可能使该关系随姿态发生小幅
变化。惯性上肢研究已表明，静态、功能性和多姿态标定的选择会显著影响最终关节角
~\cite{seel2014imu,picerno2019upper,zabat2019multiple,adjel2025influence}。本文具体的多姿态
安装级标定方法和质量门控将在第4章介绍。

\subsection{时间同步、连续性与实时输出}
\label{subsec:ch2_temporal_issue}

双肢段测量不仅要求每个传感器独立有效，还要求两个观测属于近似相同的时刻。无线接收器的
延迟抖动、相机帧率偏差或丢帧会把时间误差转化为动态姿态误差。在快速动作中，即使每个样本的
静态精度较高，几十毫秒的错位也会显著改变相对旋转。因此，实时系统应同时检查硬件时间戳、
到达时间、样本年龄、配对偏差和输出连续性。对测量系统的评价也应区分静态准确度、动态跟随、
数据有效率和端到端延迟。

\section{研究现状总结与本文测量方案}
\label{sec:ch2_gap_positioning}

现有方法不存在同时满足高空间精度、无遮挡、无限工作空间、低成本、低延迟和零安装要求的
单一方案。针对本文的近距离左上肢实时控制场景，可以得到以下结论。

\begin{enumerate}[label=\arabic*)]
  \item 标记式光学动捕适合高质量实验室参考，但场地、遮挡和成本不利于日常实时控制；
  \item 普通相机便于记录整体动作，但深度和肢段轴向旋转估计较弱，更适合作为独立评价；
  \item IMU/MIMU 便携且工作范围大，但存在漂移、动态加速度或磁场相关误差；
  \item 电磁跟踪能够在无遮挡条件下直接测得两个肢段的六自由度位姿，适合作为本文在线输入，
  但必须限制工作空间、监测金属干扰，并处理每次佩戴后的安装关系；
  \item 双手机与刚性标记能够以不同物理原理形成低成本视觉参考，用于验证电磁估计的方向、
  肘角和动态趋势，但必须报告其自身标定与同步质量。
\end{enumerate}

基于上述分析，本文采用一个固定于躯干的电磁发射器和两个分别安装于大臂、前臂的接收器，
把电磁位姿作为实时控制链的主输入；通过安装级标定将传感器状态转换为大臂方向、前臂方向及
五自由度人体状态；使用双手机刚性标记系统作为离线独立参考。该选择并不意味着电磁方法在
所有指标上优于其他技术，而是其无遮挡、直接六自由度输出和低在线计算量与本文应用边界更匹配。

\section{本章小结}
\label{sec:ch2_summary}

本章围绕人体上肢运动测量综述了标记式光学动捕、无标记视觉、惯性/磁惯性测量、电磁跟踪和
多模态视觉参考。不同技术输出的表示层级和误差来源不同，不能把关键点精度、传感器姿态精度
和人体关节角精度混为一谈。对本文采用的双电磁接收器方案而言，主要优势是无需光学视线即可
实时输出肢段六自由度位姿；主要问题则是有限工作空间、金属干扰、双接收器同步以及重新佩戴后
的传感器--肢段关系。基于这些结论，后续章节将依次给出系统设计、安装级标定与五自由度估计、
机器人动作重定向与安全执行，以及分层实验评价方法。

